%%%%%%%%%%%%%%%%%%%%%%%%%%%%%%%%%%%%%%%%%%%%%%%%%%%%%%%%%%%%%%%%%%%%%%%%
\chapter{Appendice}
%%%%%%%%%%%%%%%%%%%%%%%%%%%%%%%%%%%%%%%%%%%%%%%%%%%%%%%%%%%%%%%%%%%%%%%%

\section{Il master method}

Introduciamo brevemente un metodo per risolvere relazioni ricorsive definite per ricorrenza, detto
\emph{master method}. I lettori più interessati possono trovare una trattazione più dettagliata sul
libro~\cite{cormen2022introduction}. Il master method si usa per risolvere ricorrenze della forma
\begin{equation}
  \label{eq:master_rec}
  T(n) = aT(\frac{n}{b}) + f(n),
\end{equation}
dove $ a \ge 1 $ e $ b>1 $ sono costanti e $ f(n) $ è una funzione asintoticamente positiva. La
ricorrenza \ref{eq:master_rec} descrive il tempo di esecuzione di un algoritmo che divide un
problema di dimensione $ n $ in $ a $ sottoproblemi, ognuno di dimensione $ \frac{n}{b} $. Gli $ a $
sottoproblemi vengono risolti ricorsivamente in tempo $ T(\frac{n}{b}) $, e la funzione $ f(n) $
rappresenta il costo di dividere il problema e ricombinare la ricorsione di ogni sottoproblema.

\begin{definition}[Notazione $ O-$grande/piccolo]
  Per delle funzioni $ f,g $ a variabile reale (o intera) nella variabile $ x $ diciamo che:
  \begin{itemize}
    \item $ f(x) = O(g(x)) $ se esiste una costante $ C > 0 $ e un valore $ x_0 $ tale che $ |f(x)|
    \le C|g(x)| $ per ogni $ x\ge x_0 $.
  \item $ f(x) = o(g(x)) $ se $ \lim_{x \to \infty } \frac{|f(x)|}{|g(x)|} = 0 $.
  \item $ f(x) = \Omega(g(x))$ se esiste una costante $ C>0 $ e un valore $ x_0  $ tale che $
    C|f(x)| \ge |g(x)| $ per ogni $ x\ge x_0 $.
  \item $ f(x) = \omega(g(x)) $ se $ \lim_{x \to \infty } \frac{|g(x)|}{|f(x)|} = 0 $.
  \item $ f(x) = \Theta(g(x)) $ se $ f(x) = O(g(x)) $ e $ f(x) = \Omega(g(x)) $.
  \end{itemize}
  
\end{definition}

\begin{theorem}[Master theorem]
  \label{th:master}
  Siano $ a\ge 1 $ e $ b>1 $ delle costanti, sia $ f(n) $ una funzione e definiamo la legge
  ricorsiva $ T : \N \to \R $ con la ricorrenza
  \begin{equation*}
    T(n) = aT(\frac{n}{b}) + f(n),
  \end{equation*}
\end{theorem}
dove con $ \frac{n}{b} $ indichiamo sia $ \left\lfloor \frac{n}{b} \right\rfloor $ o $ \left\lceil
\frac{n}{b} \right\rceil  $. Allora $ T(n) $ assume i seguenti comportamenti asintotici: 
\begin{enumerate}
  \item Se $ f(n) = O(n^{\log_b a-\varepsilon}) $ per qualche $ \varepsilon >0 $, allora $ T(n) =
    \Theta(n^{\log_b a}) $.
  \item Se $ f(n) = \Theta(n^{\log_b a}) $ allora $ T(n) = \Theta(n^{\log_b a \log n}) $.
  \item If $ f(n) = \Omega(n^{\log_b a + \varepsilon}) $ per qualche costante $ \varepsilon >0 $, e
    se $ af(\frac{n}{b}) < cf(n) $ per qualche costante $ c<1$ e ogni $ n $ sufficientemente grande,
    allora $ T(n) = \Theta(f(n)) $.
\end{enumerate}

\begin{example}
  Nel caso dell'algoritmo naive la ricorrenza ha la forma $ a = 8 $, $ b = 2 $ e $ f(n) =
  \Theta(n^2) $, ovvero
  \begin{equation*}
    T(n) = 8T(\frac{n}{2}) + \Theta(n^2).
  \end{equation*}
  Siccome $ n^{\log_b a} = n^{\log_2 8} = n^3 $, che è polinomialmente più grande di $ f(n) $, ovvero $
  f(n) = O(n^{3-\varepsilon}) $ per $ \varepsilon = 1 $, dal caso 1 di \ref{th:master} segue che $
  T(n) = \Theta(n^3) $.
  Nel caso di Strassen, si ripetono i calcoli con $ a = 7 $, $ b=2 $ e $ f(n) = \Theta(n^2) $ e si
  ottiene la relazione $ T(n) = 7T(\frac{n}{2}) + \Theta(n^2)$ e analogamente $ T(n) =
  \Theta(n^{\log_2 7}) $.
\end{example}


